% Generated semantic source for AI Almanach v3.
% Edit v3 directly after initial generation; do not overwrite
% editorial changes by rerunning the converter casually.

% ==== Introduction =============================================================================
    \chapter{Introduction to AI and Data Analytics}
\chapterlead{How do intelligent systems turn goals and data into defensible decisions?}{Frame}{Represent}{Learn}{Evaluate}{orientation}

\label{ch:introduction}

\begin{learningobjectives}
After working through this chapter you should be able to:
\begin{itemize}
  \item place the main AI paradigms --- symbolic, statistical, and learned ---
        in the order they appeared and say what each one could not do;
  \item explain what caused both AI winters and recognise the pattern when it
        recurs;
  \item distinguish narrow from general AI precisely enough to spot the
        conflation in a product claim;
  \item name the four ingredients that made deep learning work after 2012, and
        identify which one was genuinely new;
  \item state, for any proposed AI application, what the non-AI baseline is and
        what evidence would justify replacing it.
\end{itemize}
\end{learningobjectives}

\begin{plainlanguage}
``Artificial intelligence'' has never named one thing. It has named a moving
target: whatever computers cannot yet do. Chess was AI until a computer won,
and then it was search. Reading handwriting was AI until it worked, and then it
was classification. This chapter follows that moving target through three
distinct technical traditions, each of which solved the previous one's blocking
problem and then met its own.
\end{plainlanguage}

\begin{engineernote}
Read the history for the failures, not the triumphs. Both AI winters followed
the same sequence: an early result that genuinely worked on small problems, an
extrapolation to problems of a different kind, funding on the strength of the
extrapolation, and a collapse when the scaling assumption proved false. The
useful skill this chapter can give you is the ability to tell a demonstration
from a system --- which is exactly the skill the current moment demands.
\end{engineernote}

    \section{Industry 4.0}

\begin{v3topic}{Information Management Cycle}
        The \emph{Information Management Cycle} describes how organisations systematically turn raw data into actionable knowledge: requirements are gathered from users, data is collected and integrated from various sources, evaluated and refined into information products, and finally deployed back to the users who initiated the need\textsuperscript{\cite{adjaouteAI2024}}.
\visualseparator
        \scalebox{\graphscale}{
            \begin{tikzpicture}[node distance=1cm and 1cm]

                % place nodes
                \coordinate (center) at (0,0) {};
                \node [block, above=of center] (information-user) {INFORMATION USER} ;
                \node [block, right=of center] (data-sources) {DATA SOURCES} ;
                \node [block, below=of center] (information-ressources) {INFORMATION RESSOURCES} ;
                \node [block, left=of center] (information-product) {INFORMATION PRODUCT} ;

                % draw edges
                \draw [line] (information-user) -| (data-sources) node[midway, above, xshift=5mm] {Requirement Analysis};
                \draw [line] (data-sources) |- (information-ressources) node[midway, below, xshift=5mm] {Integration};
                \draw [line] (information-ressources) -| (information-product) node[midway, below, xshift=-5mm] {Evaluation};
                \draw [line] (information-product) |- (information-user) node[midway, above, xshift=-5mm] {Deployment};

            \end{tikzpicture}
        }
        \captionof{figure}{The information management cycle. Note that it closes:
the users who stated the requirement are the ones who receive the product, so
a system that never returns to them has no way to learn it answered the wrong
question.}
        \label{fig:intro-information-cycle}

\end{v3topic}
\begin{v3topic}{Data Analytics, Big Data and AI Application Areas}
        \begin{multicols}{2}
            \begin{itemize}
                \item Optimize Quality Assurance
                \item Predictive Maintenance
                \item Reduce Process Runtimes
                \item Batch-size 1 Production
                \item Decision Support
                \item Minimize Ramp-up
                \item Minimize Process Error
            \end{itemize}
        \end{multicols}

\end{v3topic}
\begin{v3topic}{Data Analytics Decision Making}
\visualseparator
        \begin{tikzpicture}[node distance=0.5cm and 0.5cm, scale=\graphscale, transform shape]
                \matrix[
                    matrix of nodes,
                    nodes={ circ, text width=1cm, minimum width=1cm },
                    column sep=1cm
                ] (m1) {
                    |(big-data)| (BIG) DATA & |(smart-data)| SMART DATA & |(decisions)| DECISIONS \\
                };

                % draw arrows
                \draw [line] (big-data) -- (smart-data);
                \draw [line] (smart-data) -- (decisions);
            \end{tikzpicture}
\captionof{figure}{The four analytics questions, in increasing order of both
value and difficulty. Each level depends on the one below it, and the jump from
predictive to prescriptive is the one that requires causal assumptions rather
than more data.}
\label{fig:intro-analytics-levels}

        \vspace{2ex}

\begin{v3topic}{}
\end{v3topic}
\begin{v3topic}{Key Milestones in AI History}
        The history of AI spans over seven decades of ambition, breakthroughs, and setbacks\textsuperscript{\cite[][p.~6]{hurbansGrokkingArtificialIntelligence2020}}.
        \factvisual{
            \textbf{Read the pattern, not only the dates.}
            Early AI alternated between behavioural tests, symbolic systems,
            and trainable models. Each advance also exposed a new limitation.

            \begin{itemize}
                \item \textbf{Test:} can behaviour appear intelligent?
                \item \textbf{Name:} can intelligence become an engineering field?
                \item \textbf{Learn:} can a machine adapt from examples?
                \item \textbf{Interact:} can language create an illusion of understanding?
                \item \textbf{Limit:} what cannot a single linear unit represent?
            \end{itemize}
        }{
            \begin{tikzpicture}[x=1cm,y=0.82cm]
                \draw[almanach line] (0,0) -- (0,4.3);
                \foreach \y/\year/\event in {
                    4/1950/Turing test,
                    3/1956/Dartmouth,
                    2/1958/Perceptron,
                    1/1966/ELIZA,
                    0/1969/XOR limit}{
                    \fill[AlmanachOchre] (0,\y) circle (2.1pt);
                    \node[anchor=east,font=\scriptsize\sffamily\bfseries,
                          text=AlmanachTeal] at (-0.18,\y) {\year};
                    \node[anchor=west,font=\scriptsize\sffamily,
                          text width=27mm,align=left] at (0.18,\y) {\event};
                }
            \end{tikzpicture}
            \visualcaption{Five early milestones show the cycle from an
              intelligence criterion to a representational limit.}
              {fig:early-ai-timeline}
        }

\end{v3topic}
\begin{v3topic}{The Perceptron and the XOR Problem}
        The \textbf{Perceptron} (Rosenblatt, 1958) is a single-layer neural network: it computes a weighted sum of inputs, adds a bias, and applies a step function\textsuperscript{\cite[][p.~13]{Kriesel2007NeuralNetworks}}.
        \begin{equation}
            y = \begin{cases} 1 & \text{if } \sum_{i} w_i x_i + b \geq 0 \\ 0 & \text{otherwise} \end{cases}
        \end{equation}
\where{$x_i$ are the input features, $w_i$ the learned weights, and $b$ the bias
that shifts the firing threshold; the output is binary.}
        A single perceptron can learn AND and OR (linearly separable), but \textbf{not XOR}. XOR requires a non-linear decision boundary --- no single line can separate the classes:
\visualseparator
        \centering
        \begin{tikzpicture}[scale=1.2]
            \draw[->] (-0.3,0) -- (1.5,0) node[right] {$x_1$};
            \draw[->] (0,-0.3) -- (0,1.5) node[above] {$x_2$};
            % XOR points
            \node[circle, fill=\maincolor!80!black, inner sep=2pt, label={below left:\footnotesize 0}] at (0,0) {};
            \node[circle, fill=red!80!black, inner sep=2pt, label={below right:\footnotesize 1}] at (1,0) {};
            \node[circle, fill=red!80!black, inner sep=2pt, label={above left:\footnotesize 1}] at (0,1) {};
            \node[circle, fill=\maincolor!80!black, inner sep=2pt, label={above right:\footnotesize 0}] at (1,1) {};
            % Show no single line works
            \draw[dashed, gray, thick] (-0.2,1.3) -- (1.3,-0.2) node[right, font=\tiny] {?};
        \end{tikzpicture}
        \captionof{figure}{XOR: no single line separates 0s from 1s. A \emph{multi-layer} network (hidden layer) solves this.}
        \label{fig:xor-nonlinearity}
\end{v3topic}
    \subsection{AI Winter}

\begin{v3topic}{AI Winter: Boom, Bust, and Revival}
        The term \emph{AI Winter} refers to periods of reduced funding and interest following waves of over-optimism.\textsuperscript{\cite{adjaouteAI2024}}

        \begin{description}
            \item[First Winter (1974--1980)] The Lighthill Report\textsuperscript{\cite{lighthillArtificialIntelligenceSurvey1973}} criticised AI progress in the UK. DARPA cut funding after early machine-translation projects failed to deliver. Hardware was insufficient for the combinatorial demands of symbolic search.
            \item[Second Winter (1987--1993)] The expert-system boom of the 1980s collapsed when maintenance costs soared and the \emph{Lisp machine} hardware market evaporated. The British and US governments scaled back funding again.
        \end{description}

        \textbf{What ended them:}
        \begin{itemize}
            \item Faster CPUs and cheaper memory made previously intractable computations feasible.
            \item The revival of connectionism (backpropagation, Rumelhart \& Hinton 1986)\textsuperscript{\cite{rumelhartLearningRepresentationsBackpropagating1986}}.
            \item Commercial applications in niche domains (speech recognition, scheduling) proved ROI.
            \item The internet produced large labelled datasets at low cost.
        \end{itemize}
\end{v3topic}

    \subsection{Notable Advances in AI}

\begin{v3topic}{Landmark Achievements (1997--present)}
        A series of high-profile results demonstrated that AI could surpass human experts in well-defined domains\textsuperscript{\cite{adjaouteAI2024}}.
\visualseparator
        \begin{tblr}{width=\linewidth,
            colspec = {c X[2,l] X[5,l]},
            hlines, vlines,
            row{1} = {font=\bfseries, bg=gray!20},
            rowsep = 1pt,
        }
            Year & Name & Description \\
            1997 & Deep Blue & IBM's chess computer defeats Garry Kasparov using alpha-beta search with hand-crafted evaluation functions --- a milestone for symbolic AI. \\
            2012 & AlexNet & Krizhevsky, Sutskever, and Hinton win ImageNet with a deep CNN on GPUs, reducing error by ${\approx}10$ pp and triggering the deep-learning era. \\
            2016 & AlphaGo & DeepMind defeats Go champion Lee Sedol using deep RL and Monte Carlo tree search. Go's ${\approx}10^{170}$ positions made brute force impossible. \\
            2017 & Transformer & Vaswani et al.\ introduce \emph{Attention is All You Need}, replacing recurrence with self-attention and enabling massively parallel training. \\
            2022 & ChatGPT & OpenAI releases a dialogue-optimised LLM (GPT-3.5), reaching 100M users in two months and bringing generative AI into mainstream discourse. \\
        \end{tblr}
\captionof{table}{Landmark results, each of which was called artificial intelligence until it worked and then was renamed after its method.}
\label{tab:intro-landmarks}
\end{v3topic}

    \section{Expert Systems}
    \subsection{Overview Over Expert Systems}

\begin{v3topic}{Expert Systems Architecture}
        An \emph{expert system} is a computer program that uses a knowledge base of human expertise to solve problems within a specific domain. Expert systems were the dominant commercial AI technology of the 1980s and represent a key milestone in the practical application of AI\textsuperscript{\cite{hurbansGrokkingArtificialIntelligence2020}}.

\topicline{Core Components}
        \begin{enumerate}
            \item \textbf{Knowledge Base:} A structured collection of domain facts (\emph{assertions}) and IF--THEN production rules, typically authored by a knowledge engineer working with domain experts.
            \item \textbf{Inference Engine:} The reasoning mechanism that applies rules to the knowledge base. Two main strategies:
            \begin{itemize}
                \item \textbf{Forward Chaining} (data-driven): Start from known facts and fire applicable rules until a goal is reached.
                \item \textbf{Backward Chaining} (goal-driven): Start from the desired conclusion and work backwards to find supporting facts.
            \end{itemize}
        \end{enumerate}

\topicline{Historical Examples}
        \begin{itemize}
            \item \textbf{MYCIN} (Stanford, 1972): Diagnoses bacterial infections and recommends antibiotics.
            \item \textbf{DENDRAL} (Stanford, 1965): Identifies chemical structures from mass-spectrometry data.
        \end{itemize}

\topicline{Limitations}
        Knowledge acquisition bottleneck, poor handling of uncertainty, brittle outside domain scope, high maintenance cost.
\end{v3topic}

    \subsection{Introduction to Prolog}

\begin{v3topic}{Prolog: Logic Programming Basics}
        \textbf{Prolog} (PROgramming in LOGic) is a declarative language based on first-order predicate logic, closely associated with expert systems and symbolic AI.

        A Prolog program consists of three building blocks.
        \factvisual{
            \begin{itemize}
                \item \textbf{Facts} are unconditional assertions:\\
                    \texttt{parent(tom, bob).}\\
                    \texttt{parent(bob, ann).}
                \item \textbf{Rules} define conditional derivations:\\
                    \texttt{grandparent(X, Z) :-}\\
                    \texttt{parent(X, Y), parent(Y, Z).}
                \item \textbf{Queries} ask the inference engine:\\
                    \texttt{?- grandparent(tom, ann).}
            \end{itemize}
        }{
            \begin{tikzpicture}[node distance=6mm]
                \node[almanach card] (q) {Query\\\texttt{grandparent(tom, ann)}};
                \node[almanach accent card,below=of q] (r) {Match rule\\find an intermediate $Y$};
                \node[almanach card,below left=6mm and -2mm of r] (f1)
                    {Fact\\\texttt{parent(tom,bob)}};
                \node[almanach card,below right=6mm and -2mm of r] (f2)
                    {Fact\\\texttt{parent(bob,ann)}};
                \node[almanach accent card,below=18mm of r] (ok)
                    {Unify $Y=\texttt{bob}$\\\texttt{true}};
                \draw[almanach arrow] (q) -- (r);
                \draw[almanach arrow] (r) -- (f1);
                \draw[almanach arrow] (r) -- (f2);
                \draw[almanach arrow] (f1) -- (ok);
                \draw[almanach arrow] (f2) -- (ok);
            \end{tikzpicture}
            \visualcaption{A query succeeds when rule variables unify with
              facts that jointly prove the goal.}{fig:prolog-proof-flow}
        }

        The inference engine uses \emph{SLD resolution} (unification + backtracking) to search for proofs. Variables (uppercase) are unified with constants during resolution.

        \textbf{Relevance to AI:} Prolog underpinned early natural-language parsers, knowledge representation, and the Japanese Fifth Generation Computer project (1982). Modern applications include constraint solving and ontology reasoning.
\end{v3topic}

    \section{Neuroscience}
    \subsection{The (Human) Brain}

\begin{v3topic}{The Human Brain: Structure and Scale}
        Biological neural networks remain the primary inspiration for artificial ones\textsuperscript{\cite[][p.~13]{Kriesel2007NeuralNetworks}}

        \textbf{Key facts:}
        \begin{itemize}
            \item The human brain contains $\approx$86 billion neurons, each connected to $\approx$7{,}000 synapses on average, yielding $\approx 10^{14}$--$10^{15}$ synaptic connections.
            \item \textbf{Neuron:} The fundamental processing unit. Receives electrochemical signals via \emph{dendrites}, integrates them in the \emph{soma}, and fires an action potential along the \emph{axon} if the threshold is exceeded.
            \item \textbf{Synapse:} The junction between two neurons. Neurotransmitters cross the synaptic cleft, either exciting or inhibiting the post-synaptic neuron.
            \item \textbf{Plasticity:} Synaptic strength changes with activity (\emph{Hebbian learning}: ``neurons that fire together, wire together''), the biological basis for learning and memory.
        \end{itemize}

        \textbf{Architecture:} The cortex is organised into specialised regions (visual cortex, motor cortex, prefrontal cortex for executive function). Massive parallelism enables real-time perception far exceeding sequential digital computation at low power ($\approx$20 W).
\end{v3topic}

\begin{v3topic}{Gross Anatomy: Cortex, Deep Nuclei, and Pathways}
        Beneath the folded sheet of cortex lies a set of deep grey-matter
        structures --- \keyterm{nuclei} --- that are not decorative relay
        stations but the organising hubs of the whole
        system\textsuperscript{\cite[][chs.~1--3]{purvesNeuroscience2018}}.

        \begin{description}
          \item[Thalamus] The obligatory gateway. Almost every sensory
            signal except smell is relayed through a thalamic nucleus before
            reaching cortex --- vision through the \emph{lateral geniculate
            nucleus} (LGN), hearing through the \emph{medial geniculate
            nucleus} (MGN). Crucially, the traffic is not one-way: cortex
            sends more fibres \emph{back} to the thalamus than it receives,
            so the ``input stage'' is under continuous top-down control.
          \item[Basal ganglia] A group of nuclei --- striatum (caudate and
            putamen), globus pallidus, subthalamic nucleus, substantia nigra
            --- that select and gate actions. Dopaminergic projections from
            the substantia nigra and ventral tegmental area carry a
            \keyterm{reward prediction error}, the biological signal that
            temporal-difference learning in \cref{ch:reinforcement-learning}
            was built to mirror.
          \item[Hippocampus] Forms new declarative memories and binds them to
            spatial and temporal context. It does not store memories
            permanently: consolidation gradually transfers them to cortex.
          \item[Amygdala] Assigns emotional and threat salience, and modulates
            how strongly an experience is encoded elsewhere.
          \item[Cerebellum] Roughly $80\,\%$ of all neurons in the brain, in a
            strikingly regular circuit. Handles timing, coordination, and
            predictive correction of movement --- and, on current evidence,
            of cognition too.
          \item[Brainstem nuclei] Locus coeruleus (noradrenaline, arousal),
            raphe nuclei (serotonin), and others broadcast diffusely and set
            the global operating regime rather than carrying specific content.
        \end{description}
\end{v3topic}

\begin{v3topic}{Sensory Centres and the Hierarchies Above Them}
        Each modality has a primary cortical area that receives thalamic
        input, and a cascade of higher areas that build progressively more
        abstract descriptions\textsuperscript{\cite[][chs.~11--13]{purvesNeuroscience2018},
        \cite[][part~V]{kandelPrinciplesNeuralScience2021}}.

        \textbf{Visual pathway.} Retina $\to$ optic nerve $\to$ LGN of the
        thalamus $\to$ primary visual cortex (V1, occipital lobe) $\to$
        V2, V4, MT $\to$ two divergent streams: a \emph{ventral} (``what'')
        stream into inferotemporal cortex for object identity, and a
        \emph{dorsal} (``where/how'') stream into parietal cortex for
        location and action guidance. Receptive fields grow and become more
        selective along the hierarchy --- from oriented edges in V1 to
        object-level selectivity in inferotemporal cortex.

        \textbf{Auditory pathway.} Cochlea $\to$ cochlear nuclei $\to$
        superior olivary complex (where inter-aural timing differences
        localise sound) $\to$ inferior colliculus $\to$ MGN of the thalamus
        $\to$ primary auditory cortex (A1, superior temporal lobe) $\to$
        belt and parabelt areas $\to$ language areas in the dominant
        hemisphere. Note that the auditory pathway performs substantial
        computation \emph{below} cortex, which the visual pathway largely
        does not.

        \textbf{Convergence.} Higher association cortex --- prefrontal,
        posterior parietal, and the temporo-parietal junction --- receives
        from several modalities at once and is where perception, memory from
        the hippocampus, and value signals from the basal ganglia are
        combined into a decision.

\begin{figure}[htbp]
\centering
\begin{tikzpicture}[
  node distance=5mm,
  sens/.style={draw=AlmanachTeal,fill=AlmanachMist,rounded corners=1mm,
    minimum height=7mm,minimum width=19mm,font=\sffamily\scriptsize,
    align=center},
  nuc/.style={sens,fill=AlmanachOchre!22,draw=AlmanachOchre},
  ctx/.style={sens,fill=AlmanachTeal!18},
  hi/.style={sens,fill=AlmanachTeal!30},
  arr/.style={-{Stealth[length=1.7mm]},draw=AlmanachInk!55},
  back/.style={-{Stealth[length=1.7mm]},draw=AlmanachRed,dashed},
  lbl/.style={font=\sffamily\scriptsize,color=AlmanachInk!70}]

  \node[sens] (eye) at (0,9mm) {retina};
  \node[nuc,right=of eye] (lgn) {LGN\\\scriptsize thalamus};
  \node[ctx,right=of lgn] (v1) {V1};
  \node[ctx,right=of v1] (v4) {V2/V4/MT};
  \node[hi,right=of v4] (it) {IT\\\scriptsize ``what''};
  \node[hi,above right=1mm and 5mm of it] (par) {parietal\\\scriptsize ``where''};
  \foreach \a/\b in {eye/lgn,lgn/v1,v1/v4,v4/it}{\draw[arr] (\a) -- (\b);}
  \draw[arr] (v4) -- (par);

  \node[sens] (coch) at (0,-11mm) {cochlea};
  \node[nuc,right=of coch] (cn) {cochlear\\\scriptsize nuclei};
  \node[nuc,right=of cn] (ic) {olive /\\\scriptsize colliculus};
  \node[nuc,right=of ic] (mgn) {MGN\\\scriptsize thalamus};
  \node[ctx,right=of mgn] (a1) {A1};
  \foreach \a/\b in {coch/cn,cn/ic,ic/mgn,mgn/a1}{\draw[arr] (\a) -- (\b);}

  \node[hi] (assoc) at (108mm,-1mm) {association\\cortex};
  \draw[arr] (it) -- (assoc);
  \draw[arr] (a1) -- (assoc);
  \draw[back] (v1.south) ++(0,-1mm) -- ++(0,-4mm) -| ([xshift=-3mm]lgn.south)
    node[near end,below,lbl,color=AlmanachRed] {top-down};
  \node[lbl,anchor=west] at (0,20mm)
    {\textbf{Visual}: thalamic relay, then a long cortical hierarchy};
  \node[lbl,anchor=west] at (0,-22mm)
    {\textbf{Auditory}: much of the work happens \emph{before} cortex};
\end{tikzpicture}
\caption{The visual and auditory pathways from receptor to association
cortex, with the thalamic nuclei that gate both. The dashed arrow is the
larger, and usually forgotten, half of the story: cortex projects back to the
thalamus more heavily than the thalamus projects forward, so sensory input is
filtered by expectation before it is ever consciously available.}
\label{fig:intro-sensory-pathways}
\end{figure}

\end{v3topic}

\begin{cautionbox}[Where the brain analogy helps and where it misleads]
Three ideas in this book genuinely descend from neuroscience: hierarchical
feature extraction in convolutional networks, borrowed from the visual
pathway; attention as selective amplification; and temporal-difference
learning, which was proposed as an algorithm and \emph{later} found to match
the firing of dopaminergic neurons remarkably well --- one of the few
genuinely bidirectional exchanges between the two fields.

Four differences matter more than the similarities. Biological neurons
communicate with discrete spikes in continuous time, not real-valued
activations. Learning is local and neuromodulated; there is no known
mechanism for propagating a global error backwards through the cortex
(\cref{ch:deep-learning} covers the weight transport problem this creates).
The brain runs on about $20\,\text{W}$, roughly a thousandth of a training
cluster. And it learns continually from single examples without catastrophic
forgetting, which no current architecture does well.

Treat the analogy as a source of hypotheses, never as evidence. ``The brain
does it this way'' has motivated both major advances and long dead ends, and
you cannot tell which in advance.
\end{cautionbox}

    \subsection{Cognitive Processes}

\begin{v3topic}{Cognitive Processes and AI Analogies}
        Cognitive neuroscience studies how the brain implements perception, attention, memory, and learning -- all of which have inspired AI architectures\textsuperscript{\cite[][p.~13]{Kriesel2007NeuralNetworks}}

        \begin{description}
            \item[Perception] Hierarchical feature extraction in sensory cortices (edges $\to$ shapes $\to$ objects) inspired deep convolutional networks.
            \item[Attention] The brain selectively amplifies relevant signals and suppresses irrelevant ones. Transformer \emph{self-attention} mechanisms were directly inspired by this selective focus.
            \item[Memory] Two key systems:
                \begin{itemize}
                    \item \emph{Working memory} (short-term, $\approx$7 items): analogous to RNN hidden states or context windows in LLMs.
                    \item \emph{Long-term memory} (declarative / procedural): analogous to trained model weights.
                \end{itemize}
            \item[Learning] Reinforcement via reward (dopaminergic pathways) directly inspired temporal difference learning and reinforcement learning algorithms.
        \end{description}

        \textbf{Key caveat:} Biological and artificial neural networks share \emph{inspiration}, not mechanism. Biological neurons are far more complex than McCulloch--Pitts units; analogies should be used cautiously.
\end{v3topic}

    \section{Modern AI Systems}
    \subsection{Recent Developments in Hard- and Software}

\begin{v3topic}{Enablers of Modern AI: Hardware and Software}
        The deep-learning revolution was enabled by the convergence of three trends: massive compute, large datasets, and open-source frameworks.\textsuperscript{\cite{adjaouteAI2024}}

        \textbf{Hardware:}
        \begin{itemize}
            \item \textbf{GPUs:} Originally designed for graphics, GPUs contain thousands of small cores suited to the massively parallel matrix operations in neural network training. NVIDIA CUDA (2007) made them programmable for general AI.
            \item \textbf{TPUs:} Google's Tensor Processing Units are ASICs optimised for 8-bit matrix multiplication, offering higher throughput per watt than GPUs for inference.
            \item \textbf{Cloud Computing:} AWS, GCP, and Azure provide on-demand access to GPU/TPU clusters, removing the capital barrier to large-scale model training.
        \end{itemize}

        \textbf{Software and Data:}
        \begin{itemize}
            \item \textbf{TensorFlow} (Google, 2015) and \textbf{PyTorch} (Meta, 2016) are the dominant open-source deep-learning frameworks, providing automatic differentiation and GPU abstraction.
            \item \textbf{Hugging Face} (2016+) democratised transformer models via a shared model hub and the \texttt{transformers} library.
            \item \textbf{Big Data:} The internet produced labelled datasets at unprecedented scale (ImageNet: 14M images; Common Crawl: petabytes of text), making data-hungry models viable.
        \end{itemize}
\end{v3topic}

    \subsection{Narrow vs General AI}

\begin{v3topic}{ANI vs AGI vs ASI}
        A fundamental taxonomy in AI distinguishes systems by the breadth of their capabilities.\textsuperscript{\cite{adjaouteAI2024}}

        \begin{description}
            \item[Artificial Narrow Intelligence (ANI)] Systems that excel at a single, well-defined task (image classification, chess, speech recognition). Every deployed AI system today is ANI. Performance can exceed humans \emph{within} the domain but collapses outside it.
            \item[Artificial General Intelligence (AGI)] A hypothetical system with human-level cognitive flexibility across arbitrary tasks -- able to learn new domains, reason by analogy, and transfer knowledge. No consensus on a timeline or even a precise definition exists.
            \item[Artificial Superintelligence (ASI)] A hypothetical system surpassing human intelligence across \emph{all} cognitive dimensions. Subject of long-term safety research (Bostrom's \emph{Superintelligence}, 2014) but speculative.
        \end{description}

        \textbf{Current state:} Large language models (GPT-4, Gemini) exhibit surprising breadth but still fail on systematic reasoning and physical-world grounding. The debate whether scaling ANI converges to AGI is unresolved.
\end{v3topic}

    \subsection{NLP and Computer Vision}

\begin{v3topic}{Natural Language Processing}
        \textbf{NLP} enables machines to understand and generate human language.\textsuperscript{\cite{adjaouteAI2024}}

        \textbf{Core tasks:}
        \begin{itemize}
            \item Tokenisation, POS tagging, parsing
            \item Named Entity Recognition (NER)
            \item Machine Translation
            \item Sentiment Analysis
            \item Question Answering, Summarisation
            \item Text Generation (LLMs)
        \end{itemize}

        \textbf{Key models:} Word2Vec (2013), BERT (2018), GPT series (2018--), T5, LLaMA.

\end{v3topic}
\begin{v3topic}{Computer Vision}
        \textbf{Computer Vision} enables machines to interpret visual information from images and video.\textsuperscript{\cite{adjaouteAI2024}}

        \textbf{Core tasks:}
        \begin{itemize}
            \item Image classification
            \item Object detection (YOLO, Faster R-CNN)
            \item Semantic / instance segmentation
            \item Face recognition
            \item Optical Character Recognition (OCR)
            \item Video analysis, pose estimation
        \end{itemize}

        \textbf{Key models:} AlexNet, VGG, ResNet, Vision Transformer (ViT), SAM (Segment Anything).
\end{v3topic}

    \section{AI Application Areas}
    \subsection{Autonomous Vehicles \& Mobility}

\begin{v3topic}{Autonomous Vehicles: SAE Levels and Technology}
        Autonomous driving is one of the most visible applications of AI, combining perception, planning, and control.\textsuperscript{\cite{adjaouteAI2024}}

        \textbf{SAE Autonomy Levels (J3016):}
        \begin{center}
        \begin{tblr}{colspec={cl}, hline{1,2,8}={solid}, rowsep=1pt}
            \textbf{Level} & \textbf{Description} \\
            0 & No automation -- driver controls everything \\
            1 & Driver assistance (e.g., adaptive cruise) \\
            2 & Partial automation (Tesla Autopilot) \\
            3 & Conditional automation -- driver on standby \\
            4 & High automation in defined geographic area \\
            5 & Full automation -- no human needed \\
        \end{tblr}
\captionof{table}{The SAE automation levels. The jump that matters is from level~2 to level~3, where responsibility for monitoring transfers from the human to the system.}
\label{tab:intro-sae-levels}
        \end{center}

        \textbf{Sensor Fusion:} Autonomous vehicles combine data from LiDAR (3D point clouds), cameras (texture/colour), radar (velocity in adverse weather), and GPS/HD maps. Deep neural networks fuse these streams into a unified scene representation.

        \textbf{Challenges:} Edge cases (rare scenarios), adversarial weather, regulatory frameworks, liability, and the ``safety validation gap'' (billions of test miles required).
\end{v3topic}

    \subsection{Personalized Medicine}

\begin{v3topic}{AI in Personalized Medicine}
        Personalized (precision) medicine tailors treatment to individual patients using their genetic, environmental, and lifestyle data.\textsuperscript{\cite{adjaouteAI2024}}

        \begin{description}
            \item[Genomics] ML models identify disease-linked variants in whole-genome sequences. AlphaFold 2 (DeepMind, 2021) predicts 3D protein structure from amino-acid sequence with near-experimental accuracy, transforming drug target discovery.
            \item[Drug Discovery] Graph neural networks model molecular interactions; generative models propose novel drug candidates. Reduces early-stage discovery time from years to months.
            \item[Diagnostic AI] Convolutional networks match radiologists in detecting cancers in mammograms, skin lesions, and retinal scans. FDA-cleared AI diagnostic tools exceed 500 (as of 2023).
            \item[Clinical Decision Support] NLP extracts structured information from electronic health records (EHRs); predictive models flag deteriorating patients in ICUs.
        \end{description}

        \textbf{Challenges:} Data privacy (GDPR/HIPAA), dataset bias (underrepresented populations), regulatory approval pathways, and interpretability requirements in clinical settings.
\end{v3topic}

    \subsection{FinTech}

\begin{v3topic}{AI in Financial Technology}
        Financial services generate high-dimensional time-series data at high velocity -- a natural domain for machine learning.\textsuperscript{\cite{adjaouteAI2024}}

        \begin{description}
            \item[Algorithmic Trading] Statistical arbitrage, momentum strategies, and reinforcement-learning agents execute trades at millisecond latency. High-frequency trading accounts for $\approx$50\% of US equity volume.
            \item[Fraud Detection] Anomaly detection models (isolation forests, autoencoders, GNNs for transaction graphs) flag suspicious patterns in real time. Key challenge: class imbalance (fraud $\approx$0.1\% of transactions).
            \item[Credit Scoring] Gradient-boosted trees (XGBoost, LightGBM) supplement or replace traditional FICO scores using alternative data (payment history, behavioural signals). Raises fairness and explainability concerns.
            \item[Robo-Advisors] Automated portfolio management services (Betterment, Wealthfront) use mean-variance optimisation combined with ML risk models to provide low-cost investment advice at scale.
        \end{description}
\end{v3topic}

    \subsection{Retail \& Industry}

\begin{v3topic}{AI in Retail and Industrial Operations}
        Retail and manufacturing were among the first sectors to realise commercial value from machine learning at scale.\textsuperscript{\cite{adjaouteAI2024}}

        \begin{description}
            \item[Recommendation Systems] Collaborative filtering and deep neural networks power product recommendations (Amazon, Netflix). Amazon attributes $\approx$35\% of revenue to its recommender. Matrix factorisation and two-tower retrieval models are standard approaches.
            \item[Supply Chain Optimisation] Demand forecasting (LSTM, Prophet), inventory optimisation, and logistics routing reduce stockouts and overstock. ML scheduling cut Google's data centre cooling costs by 40\%.
            \item[Quality Control] Computer vision systems inspect products on production lines at speeds and resolutions impossible for humans. Defect detection rates $>$99\% are achievable for well-defined defect classes.
            \item[Predictive Maintenance] Sensor data (vibration, temperature, current) fed to anomaly-detection models predicts machine failures before they occur, shifting from scheduled to condition-based maintenance and reducing downtime by 20--50\%.
        \end{description}
\end{v3topic}

    \section{Further Reading}
    \begin{itemize}
        \item Chowdhary, K. R. (2020). Fundamentals of Artificial Intelligence. Springer India.
        \item Russell, S., \& Norvig, P. (2022). Artificial intelligence. A modern approach (4th ed.). Pearson Education.
        \item Ward, J. (2020). The student's guide to cognitive neuroscience (4th ed.). Taylor \& Francis Group.
    \end{itemize}
\end{v3topic}

%
\section{Deep Dive: The Foundation-Model Lifecycle}
\begin{v3topic}{One Model, Many System Layers}
A foundation model is only one component in a lifecycle that includes data
rights, pretraining, post-training, evaluation, retrieval, tools, serving,
monitoring, and incident response.
The same weights can behave very differently under another prompt template,
retrieval corpus, decoding policy, tool permission, or user workflow.

\end{v3topic}
\begin{v3topic}{Capability Is Not a Deployment Decision}
Translate a proposed capability into a user, decision, error cost, operating
boundary, baseline, acceptance test, fallback, and owner.
Multimodal input or fluent generation expands the interface; it does not remove
the need to define what counts as correct and safe.

\end{v3topic}
\begin{v3topic}{Mini Design Review}
For one AI application, draw
\emph{data $\rightarrow$ model $\rightarrow$ system $\rightarrow$ human
$\rightarrow$ outcome}.
At every arrow, name what can be lost, distorted, attacked, or misunderstood.
Then choose one metric and one qualitative investigation that could reveal the
failure before deployment.
\end{v3topic}


\begin{cautionbox}[Reading an AI claim]
Marketing language and engineering language have drifted far apart. Four
questions separate them, and they work on any claim in any domain:
\begin{enumerate}
  \item \textbf{What is the baseline?} ``$94\,\%$ accurate'' means nothing
        without the number a trivial rule achieves
        (\cref{ch:machine-learning}).
  \item \textbf{Measured on what?} A benchmark score is a statement about that
        benchmark. Ask whether the test data could have appeared in training.
  \item \textbf{Under which conditions?} Latency, hardware, input distribution,
        and failure rate at the operating threshold --- not the best case.
  \item \textbf{What happens when it is wrong?} A system with no defined
        behaviour for its own errors is a demonstration, not a product.
\end{enumerate}
A claim that survives all four is worth investigating. Most do not survive the
first.
\end{cautionbox}

\begin{checkpoint}
Pick any AI product you have seen advertised this month. Answer the four
questions above from publicly available information alone. Note which ones you
could not answer --- that gap is the actual finding.
\end{checkpoint}

\chapterreview{Frame}{Represent}{Learn}{Evaluate}
